\section{Negative character lines on the arithmetic host}
\label{sec:negative-host-lines}

% BEGIN VHN1
\paragraph{VHN1. A finite-index primitive on the original coefficient line.}
\label{claim:VHN1}
Let $E$ be a field of characteristic zero, let $D$ act on the line
$E$ by a character $\theta$, and let $z(gh)=z(g)+\theta(g)z(h)$.
Use the convention $\delta m(g)=(\theta(g)-1)m$. If
$z|_U=\delta m$ for a subgroup $U\subseteq D$ of finite index,
then $z|_D=\delta m'$ for some $m'\in E$. Indeed use the affine
action $g*x=\theta(g)x-z(g)$. The subgroup $U$ fixes $m$, so the
points $g*m$ depend only on the left cosets $gU$. Each $d\in D$
permutes this finite list and its affine map preserves their average.
Their average $m'$ is therefore fixed by $D$. This applies even if
the residue characteristic divides $[D:U]$. If $z$ and $\theta$
are integral in a DVR $O$ with fraction field $E$, clearing the
denominator of $m'$ yields $b z=\delta(bm')$ with $0\ne b\in O$
and $bm'\in O$. For any ambient group element $h$, a primitive
on $D$ gives the primitive $h*m'$ on $hDh^{-1}$ by the same action.
% END VHN1

% BEGIN VHN2
\paragraph{VHN2. Local primitives away from the residue characteristic.}
\label{claim:VHN2}
Fix a rational prime $p\ne\ell$, a valuation above $p$ in the
source algebraic closure, and its absolute decomposition group $D$.
Let $O$ be the coefficient DVR, finite over $\mathbb Z_\ell$,
$E=\operatorname{Frac}(O)$, $h\ge1$, and $\beta:D\to O^\times$
a finite character with open kernel. For the character
$\theta=\beta\chi_\ell^{-h}$ every adically continuous integral
one-cocycle has a primitive in $E$ on the whole of $D$.

To see this, pass to the open normal subgroup $U=\ker\beta$.
Its local field is a finite extension of $\mathbb Q_p$.
Write $q$ for the size of its residue field and $I_U$ for inertia.
On $I_U$ the cyclotomic character is trivial since every root of
unity of $\ell$-power order is unramified at $p$. Thus $z|_{I_U}$
is a continuous additive homomorphism to $O$. The additive group
of $O$ is pro-$\ell$, by its finite $\ell$-primary reductions.
As in LD1, wild inertia is pro-$p$ and the tame quotient of inertia
is $\prod_{r\ne p}\mathbb Z_r(1)$. Consequently $z|_{I_U}$
factors through the single factor $\mathbb Z_\ell(1)$.

Let $F_q\in U$ lift \emph{arithmetic} Frobenius. It acts on this
factor by multiplication by $q$. Explicitly, using compatible
$\ell^a$th roots $v_a$ of a local uniformizer, the tame coordinate
of $t\in I_U$ is $t(v_a)/v_a\in\mu_{\ell^a}$.
Changing the chosen root does not affect this coordinate, since
inertia fixes $\mu_{\ell^a}$. Thus conjugating $t$ by $F_q$
raises the coordinate to its $q$th power. On the coefficient
line the same Frobenius acts by $q^{-h}$. The cocycle equation
and triviality of the inertia action give
\[
     q z(t)=z(F_q tF_q^{-1})=q^{-h}z(t).
\]
Since $q>1$, $h\ge1$ and $E$ has characteristic zero, the two
scalars differ. The inertia restriction vanishes. Therefore $z|_U$
descends continuously to $U/I_U=\widehat{\mathbb Z}$. Put
$m=z(F_q)/(q^{-h}-1)\in E$. The cocycles $z$ and $\delta m$
agree on $F_q$, hence on its integer powers and by continuity on
this full quotient. Apply VHN1 to obtain an $E$-primitive on $D$.
In particular this covers primes where the original finite
character is ramified and every allowed prime in the host's $S_b$.
The argument applies to each valuation separately.
% END VHN2

% BEGIN VHN3
\paragraph{VHN3. The original negative-weight line vanishing.}
\label{claim:VHN3}
Take the assumptions of the source negative-line implication:
$\mathrm{LemmaK}(C_K)$ for the given real $C_K$, published
Poitou--Tate duality, local finiteness and local Tate duality
for both source coefficient interfaces, and local Euler
characteristic. Work on the specified algebraic closure
$\Omega/\mathbb Q$ with prime $\ell$, coefficient ring
$O=c.O\subset\overline{\mathbb Q}_\ell$, nonzero modulus $N$,
$e\ge1$, Dirichlet character $\psi$ and the original
$S,(A_v),F$ satisfying $\mathrm{IsHostFrame}$. Let
$\pi:G_{\mathbb Q}\twoheadrightarrow\Gamma=G_S$ be its quotient.
For \emph{every} $j=-n<0$ and $s\in\mathbb N$, write
$\theta=\Psi^{-n}\eta^s$. Then the source's
$\mathrm{LineVan}(F,j,\theta)$ holds. More explicitly, suppose
$z:\Gamma\to O$ is an adically continuous $\theta$-cocycle
vanishing on the prescribed $N_v$ outside $T_b$. At the
distinguished prime suppose for one $y\in O$ that
\[
 z(\mathrm{loc}_{v_\ell}d)
  =(\theta(\mathrm{loc}_{v_\ell}d)-1)y
   \qquad(d\in D_{v_\ell},\ \eta(\mathrm{loc}_{v_\ell}d)=1).
\]
There exist $0\ne b\in O$ and $m\in O$ with
$b z(g)=(\theta(g)-1)m$ for all $g\in\Gamma$.

Indeed, pull $z$ back by $\pi$ and denote the result by $Z$.
It is continuous into $E=\operatorname{Frac}(O)$: local constancy
at all powers is continuity into $O$ in its DVR topology, which
is the restriction of the finite extension topology on $E$ (RV1).
The frame gives $\Psi\pi=\chi_\ell^e\psi_{\rm Gal}$ and
$\eta\pi=\eta_4$. Set
\[
  \psi'=\psi_{\rm Gal}^{\,n}\eta_4^{-s},\qquad
             (\theta\pi)\,\psi'\chi_\ell^{en}=1 .
\]
The character $\psi'$ has open kernel: it factors through the
joint actions on $\mu_N$ and $\mu_4$. Its values, the cocycle
values and the character action lie in this same finite $E$.
There is no primeness restriction on the order of $\psi'$.
At every prime $p\ne\ell$, VHN2 applied with $h=en$ and
$\beta=(\psi')^{-1}$ gives an $E$-primitive on the whole
decomposition group, hence also on inertia. This proves the
local condition even where the native line test made no
assumption on the full inertia.

At $\ell$ the displayed hypothesis provides the primitive $y$ on
the kernel of the sign in $D_{v_\ell}$. This has index at most
two, so VHN1 gives the primitive on the full group with values
in $E$. The host's selected $A_{v_\ell}$ lies over the unique
prime of $\mathbb Q$ above $\ell$. All other prolongations there
are conjugate; VHN1 transports the primitive to each one. For
the trivial valuation the inertia condition is empty. These are
all inertia groups and all decomposition groups over $\ell$
of the original absolute local test.

We can now apply RV6 (the global vector conclusion of RV7) with
finite coefficient field $E$, finite character $\psi'$ and positive
exponent $en$. Only finiteness of the strict torsion group is
needed in RV5--6. In the proof of RV2--3 this uses finiteness
of the integral $H^2$ from the supplied $\mathrm{LemmaK}(C_K)$;
its numerical bound is finite for the given $C_K$. Thus this
application has no dependence on a particular numerical exponent
in that bound. We obtain $m_E\in E$ satisfying
$Z(g)=(\theta(\pi g)-1)m_E$ for every $g\in G_{\mathbb Q}$.
Surjectivity of $\pi$ gives this equality on $\Gamma$. Choose
one denominator $0\ne b\in O$ with $m=bm_E\in O$. The asserted
identity follows for every $g$, with one scalar and one integral
primitive for this whole line.
% END VHN3
